\section{Discussion}\label{sec:discussion}

The paper makes one claim in two registers. In random-matrix terms: the worst conditioning of a chain
of coordinate-optimized Gaussian submatrices decays at the sharp entropy rate $(k+1)H(\alpha)/\dF$, and
the optimized operator norm moves only to a bracket set by the retained-mass constants
$\kappa,\kcheck$. In network terms: a wide Gaussian network is fragile at the $1/\sqrt n$ scale at every
input, robust at exactly the reciprocal scale, and un-certifiable pattern-by-pattern below the same
entropy rate --- and all three are the same law read through the dictionary that gates are coordinate
selections and the Jacobian is a submatrix product.

Three features distinguish the account from its predecessors. It needs no oddness or bias-free
hypothesis: the small output comes either from a random readout, which is symmetry-free, or from
balance, which covers trained classifiers, and the gradient facts come from submatrix conditioning,
which does not see the sign structure. It is two-sided: the same bracket that produces the attack
produces a matching certificate, so the scale is pinned rather than bounded on one side. And it is
uniform over inputs, because the arrangement bound makes the union over realized gate patterns
affordable in the wide regime. The cost is the wide regime itself, $dL\log\bar n=o(\min_\ell n_\ell)$,
and the Gaussian (or subgaussian, or Gaussian-plus-bounded-deviation) weights; the fully trained,
proportional-width regime is reached only through the perturbation-stability theorem, which keeps the
scale but not the constants.

Several questions remain. The product rate is proved for equal densities and square bottlenecks; the
unequal-density cascade of \Cref{rem:unequal} should give the general rate, but the sharp constant for a
genuine rectangular bottleneck chain is open, as is the corresponding operator-norm edge for unequal
densities \cite{shmaloProducts}. The transformer theorem covers the simplified, standard-temperature
architecture at random init in the wide regime; the fully standard transformer, with biases, positional
embeddings, affine normalization, and GELU, breaks the balance that supplies the small output, and the
random-readout route does not obviously apply because the readout is a learned class token. The
certification gap is stated for pattern-uniform certificates; whether a certificate that adapts to the
realized pattern but is stable under a single gate flip can beat the $1/\sqrt n$ scale is the precise form
of the ``can we certify robustness'' question, and \Cref{rem:cert-interp} suggests the answer is no
without controlling the flip, which is what the adversary controls. Finally, the weight-tied dynamical
limit shows the eigenvalue radius, not the operator norm, governs the sustained gain; making the
infinite-time contraction unconditional requires the multi-time sign-decoupling estimate identified in
\cite{relumap}, which is the one genuinely open problem in the reused-weight theory.

The unifying point is methodological. Adversarial fragility has been explained by feature compression
\cite{fawzi18}, by computational hardness \cite{bubeckSS19}, and by isoperimetry
\cite{shafahi19,advExt}; the present account locates it in the spectral theory of the network's own
Jacobian, and shows that the same spectral object --- a product of Gaussian submatrix blocks, optimized
over supports --- controls the attack, the certificate, and the gap between them, quantitatively and with
matching constants. The random-matrix theorems were the missing ingredient; with them the geometry of
adversarial examples in deep networks is, in the wide Gaussian regime, a corollary of the conditioning
of coordinate-optimized Gaussian maps.
